\chapter{The USSR and the East (1923--1927)}

The non-European countries occupied only a peripheral place in Marx's thought, and were neglected by the First and Second Internationals. When Lenin, in his famous work published in 1916, diagnosed ``imperialism'' as the highest, and last, phase of capitalism, he was more concerned with its implications for the imperialist countries than for their colonial subjects. The proclamations to the peoples of Asia during the first year of the revolution were for the most part incitements to revolt against foreign, and in particular British, rule; and the founding congress of Comintern in March 1919 included in its manifesto an appeal to the ``colonial slaves of Asia and Africa''. It was the second congress in June 1920 which first sought to lay down a policy for what were called ``the colonial and semi-colonial countries''. Theses drafted by Lenin called for ``a close alliance of all national and colonial liberation movements with Soviet Russia''. Whether the national movements with which this alliance would be struck would be bourgeois-democratic or proletarian-communist depended on the stage of development of the country concerned. In backward countries communists must be prepared to support every ``national-revolutionary'' movement of liberation, even of a bourgeois-democratic character. It was a commonsense solution, which continued to present many practical problems. 
% 注：修复了 "pub- lished" 和 "incite- ments" 等跨行断字。

After the congress, Comintern took its first major initiative in eastern affairs by convening a ``congress of peoples of the east'' at Baku, which mustered nearly 2000 delegates, most of them from Central Asia, and predominantly Muslim. Throughout this area it was not difficult to depict British imperialism as the major enemy; and this was the main theme of orators at the congress. But embarrassments arose both from the religious susceptibilities of many Muslim delegates, and from the presence of Enver, a leader of the Young Turk nationalist revolution of 1908, who was widely held responsible for the massacres of Armenians, and whose socialist or democratic credentials were conspicuously defective. The congress had no sequel, and yielded no lasting results. A year later a similar congress of Far Eastern peoples was projected at Irkutsk. The plan fell through, and the congress was eventually held in Moscow in January 1922. But by this time enthusiasm had waned, and it proved less impressive than its predecessor at Baku. In the Far East Japan was the country where industrialization had progressed far on a western model, which possessed a numerous proletariat, and therefore seemed to offer the most promising prospect of revolution. But no Japanese delegates appeared at the founding congress of Comintern; and capitalist Japan was even more impenetrable to the impact of communism than the capitalist countries of the west. It was China, where a growing national movement for liberation from imperialist domination was directed against the ``unequal treaties'' and foreign settlements in the ``treaty ports'', which proved the most fruitful field for communist propaganda and Soviet diplomacy. 
% 注：修正了 "diplo- macy" 断字。

Lenin, in an article of 1912 inspired by the Chinese revolution of that year, had declared that ``the hundreds of millions of toilers in Asia have a reliable ally in the proletariat of all civilized countries'', and predicted that the victory of the proletariat ``will free the peoples of Europe and the peoples of Asia''; and he described Sun Yat-sen, the Chinese nationalist leader, as a narodnik with a ``revolutionary-democratic core'' in his programme. When in 1918 Sun Yat-sen set up a dissident nationalist government in Canton, which made him the recognized head of the national movement, mutual sympathy between the two revolutionary centres was demonstrated in an exchange of letters and telegrams between Sun and Chicherin. Early in the nineteen-twenties Soviet diplomacy first became active in China. Japanese troops had lingered in Siberia long after the other Powers active in the civil war withdrew their units. But by 1921, under American pressure, they were on the way out. Soviet forces moved gradually eastwards, expelled a White Russian army which had occupied Outer Mongolia, and in November 1921 proclaimed a Mongolian People's Republic under Soviet patronage and control. In the summer of 1922 Joffe was despatched from Moscow in an attempt to clear up relations with the shadowy and largely impotent Chinese Government in Peking. The attempt failed. But in January 1923 Joffe had a meeting in Shanghai with Sun Yat-sen, who had been recently driven out of Canton. It was a moment when the principles of the united front and of cooperation with national movements to resist imperialism were firmly implanted in Soviet policy. A joint declaration signed at the end of the meeting recorded Joffe's acceptance of Sun Yat-sen's view that ``neither the communist order nor the Soviet system can actually be introduced into China, because there do not exist there the conditions necessary for the successful establishment of either communism or Sovietism''. But it was agreed that ``China's paramount and most pressing problem is to achieve national unification and attain full national independence''; and Joffe gave an assurance that China could count on the warmest sympathy and support from Russia in this task. 
% 注：修复了 "Ameri- can" 和 "in- troduced" 的跨行断字。

Two months later Sun Yat-sen regained power in Canton; and the agreement with Joffe was the starting-point of a long and fruitful period of Soviet cooperation with Sun and his party, the Kuomintang. In the autumn of 1923 Chiang Kai-shek, one of Sun's lieutenants, was sent to Moscow to negotiate for the supply of arms and equipment; and Borodin, a Russian, American-born, English-speaking communist, arrived in Canton to act as adviser to Sun. During the next year Borodin succeeded in cementing a close alliance between himself and Sun, and between the Soviet Government and Kuomintang, the common aim being to liberate China from the domination of the imperialist Powers---Britain, Japan and the United States. Sun, since his return to Canton, had established there a nationalist government, which planned one day to launch a military ``northern expedition'' to reunite China and drive out the privileged foreign intruders. Military supplies reached Canton in modest but increasing volume from the USSR; Soviet military advisers helped to build up the Canton army, and to equip and staff a new Military Academy. Sun, with guidance from Borodin, tightened up the loose organization of Kuomintang. The Chinese Communist Party (CCP), founded in 1921, had at this time scarcely more than a thousand members, mainly Marxist intellectuals. Before Borodin's arrival, and apparently at the instigation of Comintern, an agreement had been reached under which members of the CCP were also to become members of Kuomintang. The model was apparently the dual status of many members of the CPGB as members of the Labour Party; and the intention was that this disciplined and devoted group should add a stiffening to the larger and looser organization. All these arrangements masked the discrepancies between Marxist doctrine and the ``three principles'' of Sun Yat-sen---``nationality'', ``democracy'' and ``people's livelihood''. This was easy so long as everything else was subordinated to the national revolution against imperialism. It was only when Borodin pressed for the inclusion of the expropriation of landlords in the Kuomintang programme that Sun stubbornly resisted, and Borodin had to give way. 
% 注：修复了 "v/ith" 为 "with" 以及去除了 "(land the intention" 中的杂乱括号。

At the end of 1924 Sun Yat-sen embarked on a journey to Japan and northern China to take stock of the situation. He fell ill on the way, and died in Peking on March 12, 1925. The succession seemed most likely to fall on Wang Ching-wei, a clever but weak man, who belonged to the Left wing of Kuomintang. Chiang Kai-shek's military capacities, and the prestige which he had acquired by his visit to Moscow, gave him a commanding position. But for the moment he revealed no political ambitions, and leaned heavily on Soviet support in building up the national army. The most sensational event of 1925 occurred on May 30 in Shanghai, when municipal police under British command fired on a demonstration of striking workers and students, killing several of them. This act provoked a general strike and mass labour disorders which lasted for two months, and spread to Canton. For the first time an effective trade union organization was formed in Shanghai under CCP leadership; and membership of the CCP leapt up in a few weeks to 10,000. One result of these first symptoms of workers' revolt in China was to sharpen mutual animosity between Britain and the USSR. Another was to encourage the growth in Kuomintang of a Right wing committed to the cause of national liberation, but hostile to social revolution. Chiang Kai-shek watched events, and manoeuvred quietly between Left and Right. 
% 注：去除了 ".Another" 中的前导点。

Concern for the national-revolutionary movement centred on Canton did not exhaust Soviet interest in China. It was northern China which abutted directly on Soviet territory. In August 1923 Karakhan had arrived in Peking as diplomatic representative to the Chinese Government, and concluded in May 1924 a treaty for the regularization of Sino-Soviet relations. The Soviet Government had already renounced the extra-territorial rights and concessions which Russia, in common with the other major Powers, had formerly enjoyed in China. The remaining bones of contention were Outer Mongolia, over which the Chinese Government still claimed sovereignty, and the Russian-built Chinese Eastern Railway (CER), which traversed the length of Manchuria on its way to Vladivostok. Outer Mongolia was recognized by the treaty as ``an integral part'' of China; but no date was fixed for the withdrawal of Soviet troops or administrators, and the USSR was determined to keep a firm hand over the Mongolian People's Republic. The CER was placed under the control of a board consisting of five Chinese and five Russian members; but the general manager of the line was appointed by the Soviet Government---an arrangement which caused much friction in the years to come. The Soviet Government was conscious of no incompatibility between the defence of its interests in north China and the promotion of the revolutionary cause in the south. But some circles in Kuomintang keenly resented the conclusion of these agreements by the USSR with sworn enemies of the nationalist movement. 
% 注：修正了 "the the Chinese" 的叠词，并去除了 "uscious"、"' in the south"、"' the conclusion" 的干扰符。

The Chinese Government in Peking which had negotiated the Sino-Soviet treaty was under the loose control of Wu Pei-fu, the warlord who had for some time dominated central China, and enjoyed British support. In the autumn of 1924 hostilities broke out between Wu and Chang Tso-lin, the warlord of Manchuria and a protege of Japan. Wu's defeat was hastened by the desertion of Feng Yu-hsiang, who held a large area in north-western China. Hitherto a subordinate of Wu, Feng now announced his sympathy with Kuomintang and the nationalist government in Canton---a change of front which may have been encouraged by subsidies and offers of support from Moscow. After the decline of Wu's authority Feng sought to assert his own control over Peking and the adjacent provinces. But this ambition was thwarted by Chang Tso-lin, who drove him out at the end of 1925. Thereafter, the hapless Peking government became the puppet of Chang. There were now only two major military forces in China: those of Chang Tso-lin in the north, and the rapidly expanding nationalist forces commanded by Chiang Kai-shek in the south. 
% 注：修复了 "1925-" 为 "1925."，并增加了冒号使 "forces in China those of..." 逻辑通顺。修正 "Feng Yii-hsiang" 为 "Feng Yu-hsiang"。

Much of central China was a prey to the disintegrating armies which had once owed allegiance to Wu Pei-fu. It was in these conditions that Chiang took, at the beginning of 1926, the momentous decision to start the long-projected ``northern expedition'' in the summer. It was not welcomed by Borodin or by the Soviet advisers. The northern expedition had been talked of continually as the ultimate objective of the military preparations, and applauded in principle. But, presented as a concrete plan for the immediate future, it inspired apprehension. Its success was uncertain, and it seemed likely to provoke intervention by the imperialist Powers. The Soviet Government was much alarmed at this time by a dispute with Chang Tso-lin over the CER, wanted no trouble elsewhere, and paid no great attention to events in Canton. Borodin left Canton in January 1926 for a visit to Peking and to Feng's headquarters; and, while he was away, a quarrel broke out between Chiang and the senior Soviet military advisers, who were tactlessly sceptical of the projected enterprise. On March 20, 1926, a trumped-up incident, arising out of the movements of a Chinese gun-boat whose commander was a communist, gave Chiang an excuse for confining several Soviet advisers to their quarters and arresting Chinese communists attached to the armed forces. The advisers were quickly released; but Chiang peremptorily demanded the departure of those who had disputed his authority. When Borodin returned to Canton at the end of April, peace had been restored and honour saved all round. The offending advisers were withdrawn. Blyukher (alias Galin), a Red Army officer who had previously served in China and was persona grata to Chiang, arrived to take charge of an enlarged team of Soviet military advisers. Everyone now accepted the imminence of the northern expedition, and Blyukher and his staff worked actively to plan and organize it. But the balance of forces had changed. Chiang was firmly in command. 

Early in July 1926 the nationalist army 70,000 strong, with a full complement of Soviet advisers, marched northwards from Canton. The campaign was a brilliant success. Not only was no resistance encountered, but large reinforcements---units from Wu Pei-fu's disbanded armies, and groups of armed peasants living on the plunder of landlords' estates---joined the march. When Chiang at the beginning of September entered Hankow, the great industrial city of central China and the former capital of Wu's dominions, his force numbered some 250,000. A few weeks later he moved eastward to set up his head-quarters at Nanchang---a first step on the road to Shanghai. In November the Kuomintang authorities left behind in Canton, together with Borodin and his staff, journeyed to Hankow, where a national-revolutionary government was proclaimed amid scenes of enthusiasm. The city was enlarged by the accretion of two adjacent industrial centres and renamed Wuhan. It was a moment of triumph, in Wuhan and in Moscow. 

Victory concealed, however, the seeds of calamity. So long as the revolutionary movement remained within its nationalist framework, and preached liberation from foreign imperialism, unity prevailed. But, when some of its sponsors began to speak of the liberation of peasants or workers from feudal or capitalist oppression, jarring notes recurred. Kuomintang was predominantly petty bourgeois. It had more small landowners than landless peasants among its members; most officers in the nationalist forces were said to own land. Nor had it any specific links with the workers, or with the trade union movement in Shanghai initiated by the events of May 30, 1925. The session of IKKI in Moscow in November 1926, which hailed the victory of the Chinese revolution, gave an uncertain lead. It looked forward to the next stage of the revolution, in which the proletariat would take the lead; and it proclaimed the importance of agrarian revolution in China. It instructed Chinese communists to remain in Kuomintang and support the national movement. The CCP was divided and hesitant. But Borodin interpreted the views of Moscow correctly when he insisted on its loyal support for Kuomintang, even if this involved postponement of the demands of workers and peasants to a more convenient season. 
% 注：去除了 "Iwas"、"'■ union"、"I I of the revolution"、"I lead"、"Hn China" 的边缘干扰符。

The crisis came through a split in Kuomintang itself. The Wuhan government, representing the Left wing of Kuomintang and strongly influenced by Borodin, combined support of the national revolution with much lip-service to the aims of social revolution. Peasant revolts were rife in Hunan, the province to the south of Wuhan; and this was the moment when Mao Tse-tung first became prominent as the champion of the peasants. In Nanchang, Chiang Kai-shek and his generals moved sharply to the Right, expressing open hostility to the communists and to the demands of unruly peasants and workers, which interfered with his nationalist ambitions. These developments were assisted by the changed attitude of the British Government, which, impressed by the sweeping success of the nationalist forces, concluded that it might be wiser to come to terms with them than to fight them. It paved the way to an agreement by returning to Chinese control the British concessions in Hankow and Kiukiang, and proposing to relax or abolish other servitudes imposed on China by the unequal treaties of the past. Chiang, who had long been impatient of Soviet tutelage, and was now strong enough to dispense with it, perceived a dazzling opportunity of realizing his ambitions with the blessing of the imperialists, whose antipathy to the communists and to their programme of social revolution ran parallel with his own. 
% 注：去除了 "'of the peasants" 和 "'antipathy" 的乱码单引号。

The full implications of the change were not at once realized. Shanghai was controlled at this time by a minor warlord, Sun Ch'uan-fang, who was clearly in a vulnerable position. In February 1927 the trade unions in Shanghai organized a workers' rising, counting on aid from Chiang Kai-shek, whom they still looked on as a liberator. Chiang made no move; and Sun easily suppressed the rising. A few weeks later Sun's forces were defeated by those of Chiang in a pitched battle outside Shanghai. Once more the workers in Shanghai rose, set up organs of local self-government, and prepared to welcome the entry of nationalist forces into the city. When Chiang eventually arrived, his disapproval of these proceedings was made manifest. Order was enforced by the troops, and organs of government disbanded. Then, on April 12, when everything was ready, Chiang let loose a large-scale organized massacre of communists and militant workers throughout the city. The CCP and the trade unions were wiped out. This time the message was unmistakable. Chiang was fiercely denounced in Wuhan and in Moscow. But these protests did not alter the fact that Chiang commanded the one effective army in central and southern China, and had won the sympathy and toleration of the foreign Powers. 
% 注：去除了 "'12" 和 "(throughout" 的边缘乱码。修复了跨行断开的句子 ("/the one effective army... [won the sympathy... ")。

A few days earlier another disaster had overtaken Soviet policy and Soviet prestige in China. On the orders of Chang Tso-lin, and with the complicity of the diplomatic corps, the Peking government conducted a raid on the Soviet Embassy. The Ambassador's residence was spared, but the outlying buildings were ransacked, employees arrested, and a mass of papers seized. The Chinese employees were summarily executed, the Soviet employees held in prison for many months awaiting trial. Large numbers of documents, some genuine, some suitably doctored, were published in various languages by way of demonstrating the communist conspiracy against the established order. Soviet protests fell on deaf ears, and diplomatic relations were broken off. These events preceded by a month the Arcos raid in London and the rupture of Anglo-Soviet relations. 
% 注：去除了 "/by way" 和 "‘ diplomatic" 的多余符号。

In the summer of 1927 Soviet fortunes in China reached their nadir. In Wuhan, the local warlord declared his independence of Chiang. But he had no more sympathy than Chiang for social revolution, and carried out a massacre of peasants in Changsha, the capital of Hunan. Borodin and the Wuhan government counted on the loyalty of Feng Yu-hsiang, who had just returned from a long and apparently enthusiastic visit to Moscow. But Feng preferred to do a deal with Chiang, as a result of which he dismissed his Soviet advisers, and debarred communists from working in his army. The CCP held a congress at Wuhan in April-May 1927, at which it claimed a membership of 55,000. But its impotence was apparent. The Wuhan government slowly disintegrated. One of its last acts was to demand the recall of Borodin. He left China at the end of July; and the last Soviet military advisers and members of other Soviet missions also took their departure. Of four years of feverish effort directed from Moscow nothing seemed to survive. Blows had been struck from which even the most optimistic observers could see small hope of recovery. During these years a vast revolutionary ferment had indeed been generated all over China. But it remained, for a long time to come, effectively crushed under the iron heel of Chiang Kai-shek. 
% 注：去除了 "I which" 和 "I hope" 这种由折痕引起的边缘干扰线。

Ambitious schemes were mooted from time to time to extend communist propaganda and influence over the Pacific area, seamen being regarded as the most promising agents for such work. A conference of Pacific transport workers (mainly seamen, though some railway workers were also represented) met in Canton in the summer of 1924, apparently under mixed communist and Kuomintang auspices. More than twenty delegates were present from north and south China, from Indonesia and from the Philippines; Japanese delegates were prevented from making the journey. The conference sent greetings to Comintern and Profintern. But its platform seems to have been anti-imperialist rather than specifically communist. Nothing further happened till the summer of 1927, when another Pacific conference was held in Wuhan. This time Lozovsky, the president of Profintern, had come from Moscow, and the conference proceeded under his masterful direction. Delegates attended from the USSR and China, from Japan, Indonesia and Korea, and from Britain, France and the United States; delegates from Australia and India failed to arrive owing to the veto of their respective governments. The conference proclaimed its support of the Chinese revolution, demanded independence for Korea, Formosa, Indonesia and the Philippines, and set up a permanent Pan-Pacific Secretariat which led a rather shadowy existence for some years in different centres, and published a periodical called the Pacific Worker. 
% 注：修复了 "Formo.sa" 中的点号误识。

Other parts of the eastern world were less open in this period to the activity either of the Soviet Government or of Comintern. Soviet relations with Japan were not unlike those with other capitalist countries. Once Japanese troops had been withdrawn from the mainland of Siberia, the most important Soviet demands were for the evacuation of northern Sakhalin and for diplomatic recognition. Both these were belatedly achieved in a treaty of January 1925. But the questions of fishery rights and of competition between the CER, which fed Vladivostok, and the Japanese South Manchurian Railway, which fed the Japanese-controlled port of Dairen, were constant causes of friction; and mutual suspicion continued to cloud relations. Initial faith in the revolutionary potential of the Japanese proletariat was not fulfilled. The Japanese police was ruthless and efficient; and the first Japanese Communist Party dissolved itself early in 1924. It was reconstituted as an illegal organization in December 1926. Some trade unions joined a dissident federation with Left or communist affiliations. But these efforts achieved little except, from time to time, to exacerbate Soviet-Japanese relations; and the party was once more virtually stamped out by wholesale arrests in 1929. 
% 注：修复了 "reconsti- tuted" 和 "com- munist" 的跨行断字。

Elsewhere there was little to record. The cause of Indian communism enjoyed little success, except among Indians living in Europe. A small communist party led a precarious existence, constantly harassed by the British authorities. Provincial workers' and peasants' parties promoted by communists showed more promise. The demands of the Indian National Congress for independence or autonomy were widely supported; and protests against the dilatory and half-hearted concessions offered by the British Government were frequent. Some strikes were said to have been fomented by communist propaganda. But the government had the situation well in hand. In Indonesia, a small communist party was reinforced by a popular Muslim nationalist organization and by an incipient trade union movement. In November 1926, apparently without prompting or support from Comintern, it staged a mass rising, which was crushed within a few days. Executions and mass deportations followed, and put an effective end to the Indonesian party for many years. The Middle East offered still fewer opportunities for Soviet diplomacy or for communist infiltration. Relations with Turkey and Persia were designed to counteract the influence of the western Powers, and particularly of Britain, in these countries, and to develop trade between them and the USSR. Occasional embarrassments in dealing with regimes fiercely repressive of all movements of the Left did not disturb the course of Soviet policy. In Egypt the national movement of revolt against British domination grew slowly, and had no affiliations with the USSR. The Arab countries, as well as Palestine, were still too firmly under western control to permit of any significant Soviet or communist activity. 
% 注：去除了 "iThe" 和 "'under" 等乱码。